\documentclass[11pt]{article}
\usepackage{preamble}
\setlength{\parskip}{4pt}

\begin{document}

\daybanner{AP Statistics \textperiodcentered\ Unit 2 \textperiodcentered\ Topic 2.7 \textperiodcentered\ B: 2026-10-29 \textperiodcentered\ E: 2026-11-18 \textperiodcentered\ TEACHER KEY}{Overlap Does Not Decide Independence}

\sectionbanner{CED TETHER}

{\small
\begin{itemize}[leftmargin=1.5em,itemsep=2pt,topsep=2pt]
  \item \textbf{Skill 3.A} --- Determine relative frequencies, proportions, or probabilities using simulation or calculations.
  \item \textbf{EU VAR-4} --- The likelihood of a random event can be quantified.
  \item \textbf{LO VAR-4.E} --- Calculate probabilities for independent events and for the union of two events. [Skill 3.A]
  \item \textbf{EK VAR-4.E.1} --- Events A and B are independent if, and only if, knowing whether event A has occurred (or will occur) does not change the probability that event B will occur.
  \item \textbf{EK VAR-4.E.2} --- If, and only if, events A and B are independent, then P(A | B) = P(A), P(B | A) = P(B), and P(A $\cap$ B) = P(A) $\cdot$ P(B).
  \item \textbf{EK VAR-4.E.3} --- The probability that event A or event B (or both) will occur is the probability of the union of A and B, denoted P(A $\cup$ B).
  \item \textbf{EK VAR-4.E.4} --- The addition rule states that P(A $\cup$ B) = P(A) + P(B) - P(A $\cap$ B).
\end{itemize}
}
\textbf{Essential question:} Which probability comparison determines independence, and how is that different from checking overlap?\par
\textbf{DOK-3 skill:} 4.B \textperiodcentered\ \textbf{FRQ pattern:} {\small\texttt{independent-\allowbreak{}vs-\allowbreak{}mutually-\allowbreak{}exclusive-\allowbreak{}decide-\allowbreak{}from-\allowbreak{}the-\allowbreak{}table}}\par
\textbf{DOK rationale:} Part (c) weighs competing claims using the day's numerical or graphical evidence and explains a limitation or consequence; the judgment requires linked reasoning beyond a calculation.\par

\frameworkphaseheader{Do Now (first take) $\rightarrow$ Explore (video + follow-along) $\rightarrow$ Exit (finish + turn in)}{1 $\rightarrow$ 3}{45}{%
\begin{itemize}[leftmargin=1.5em,itemsep=2pt,topsep=2pt]
  \item Show the board at the bell and collect a one-sentence first take without confirming a claim.
  \item Run the named video follow-along, then allow about 10 minutes for parts (a)--(c).
  \item Ask for evidence linked to a claim. Collect the sheet and hand-score part (c) with E/P/I.
\end{itemize}
}{%
\begin{itemize}[leftmargin=1.5em,itemsep=2pt,topsep=2pt]
  \item Commit to an independence decision and cite one table feature.
  \item Complete the video follow-along about independence, intersections, and unions.
  \item Finish (a)--(c), revise the first take in the exit line, and turn in the sheet.
\end{itemize}
}{%
\begin{itemize}[leftmargin=1.5em,itemsep=2pt,topsep=2pt]
  \item Which number or visible feature supports your claim?
  \item Why does the other claim go beyond that evidence?
  \item What would you tell someone who chose the other claim?
\end{itemize}
}{Circulate and ask students to connect their choice to evidence. Score part (c) using the three required reasoning elements; do not require dropped extension work.}

\begin{callout}{\IconWarn\ WATCH FOR}{calloutred}
\begin{itemize}[leftmargin=1.5em,itemsep=2pt,topsep=2pt]
  \item A choice with copied numbers but no explanation of how they support it.
  \item Treating a model or average as a guaranteed individual outcome.
\end{itemize}
\end{callout}

\sectionbanner{SCORING PART (c) --- E / P / I}

\textbf{Expected elements}
\begin{itemize}[leftmargin=1.5em,itemsep=2pt,topsep=2pt]
  \item \textbf{reasoning-1} (required) --- Rejects independence using 0.70 versus 0.50
  \item \textbf{reasoning-2} (required) --- Uses the overlap of 84 to reject mutual exclusivity
  \item \textbf{reasoning-3} (required) --- Explains why overlap and probability change answer different questions, supporting Marcus
\end{itemize}

{\small\renewcommand{\arraystretch}{1.25}
\noindent\begin{tabular}{|p{0.5in}|p{5.9in}|}
\hline
\textbf{E} & Includes all three required reasoning elements above, with correct contextual evidence. \\ \hline
\textbf{P} & Includes at least one substantive correct reasoning link above, but omits or misstates another required element. \\ \hline
\textbf{I} & Includes no substantive correct reasoning link above; an unsupported choice or copied number alone is insufficient. \\ \hline
\end{tabular}}

\textbf{Common mistakes}
\begin{itemize}[leftmargin=1.5em,itemsep=2pt,topsep=2pt]
  \item Comparing $P(I\mid P)$ with $P(P)$ instead of with $P(I)$
  \item Saying events must be independent whenever both can occur
  \item Calling owns a pet and plays an instrument mutually exclusive despite the count 84
  \item Using $P(P)P(I)$ for the actual intersection before independence has been established
\end{itemize}

\clearpage
\focusbanner{TODAY'S PROBLEM --- annotated}

Suppose a school surveys 200 students about whether they own at least one pet and whether they play a musical instrument. Let $P$ be the event that a randomly selected surveyed student owns a pet and $I$ the event that the student plays an instrument.\par\smallskip \textbf{Aaliyah:} \emph{The events overlap because 84 students do both. Since they are not mutually exclusive, they must be independent.}\par \textbf{Marcus:} \emph{Overlap does not decide independence; the relevant probabilities must be compared.}\par
\begin{center}
\textbf{Pet ownership and instrument playing}\par\smallskip
\begin{tabular}{|l|c|c|c|}
\hline
\textbf{Pet?} & \textbf{Plays} & \textbf{Does not} & \textbf{Total} \\ \hline
\textbf{Yes} & 84 & 36 & 120 \\ \hline
\textbf{No} & 16 & 64 & 80 \\ \hline
\textbf{Total} & 100 & 100 & 200 \\ \hline
\end{tabular}
\end{center}
\begin{firsttakebox}{FIRST TAKE (5 min)}
Does owning a pet seem to change the chance that a student plays an instrument? Give one table detail in one sentence.\par
\answer{Not graded. Aaliyah reaches the correct dependence decision from an invalid overlap rule; preserve that tension.}
\end{firsttakebox}

\textit{Rules callout ``INDEPENDENCE, INTERSECTION, AND UNION (keep on the page)'' is printed on the student sheet.}\par\medskip

\dokbadge{1}~\textbf{(a)} State the counts for $P$, $I$, and $P\cap I$.\par
\answer{There are 120 students in $P$, 100 in $I$, and 84 in $P\cap I$.}

\dokbadge{2}~\textbf{(b)} Find $P(I)$, $P(I\mid P)$, and $P(P\cup I)$. Show the fractions or setup.\par
\answer{$P(I)=100/200=0.50$; $P(I\mid P)=84/120=0.70$; $P(P\cup I)=(120+100-84)/200=0.68$.}

\dokbadge{3}~\textbf{(c)} Whose reasoning is supported? Decide independence using the two instrument-playing probabilities. Explain why the 84 students in both groups settle mutual exclusivity but do not settle independence. Write 3--4 sentences.\par
\answer{Marcus's reasoning is supported: the events are not independent because $P(I\mid P)=0.70$ differs from $P(I)=0.50$. Knowing pet ownership changes the chance of instrument playing. The 84 students who do both show the events are not mutually exclusive, but overlap alone does not tell us whether one event changes the probability of the other.}

\begin{summaryexitbox}{EXIT}
One thing the video changed about my first take: \hrulefill
\end{summaryexitbox}

\end{document}
